% !TeX root = ../../tesi.tex

The implemented framework follows a shared offline--online logic across the two industrial sites, while adapting the final processing stage to the different monitoring objectives of raw dough and baked bread. In both cases, the first phase consists of autonomous image acquisition on an embedded edge device. The collected data are then transferred to an offline workstation, automatically labelled, and used to train a DEIMv2 detector. Once training is completed, the detector is redeployed on the edge hardware and becomes the entry point of the online inspection pipeline.\cite{labby_thesis_repo_2026,labby_kinect_acquisition_2026} The overall architecture was designed in a \textbf{modular} way, so that different detection, reconstruction, and downstream-analysis modules could be tested, replaced, or extended without rewriting the entire workflow.\cite{labby_thesis_repo_2026}

At \textbf{Site~1}, the acquisition stage relies on synchronized RGB-D frames from the Azure Kinect.\cite{microsoft_azure_nodate} After online detection, each dough ball is associated with aligned depth information and converted into a local 3D representation for volumetric analysis. The adopted workflow includes an initial color-clustering stage, subsequent foreground segmentation and refinement with respect to the reference table model, 3D completion of the observed geometry, and volume estimation through multiple geometric methods.\cite{labby_thesis_repo_2026}

At \textbf{Site~2}, the acquisition stage is based on RGB images from the Raspberry Pi Camera Module~3 Wide.\cite{raspberrypi_camera_module_3_2026} After online detection, the localized bread loaves are cropped and forwarded to an anomaly-detection stage implemented with EfficientAD through anomalib.\cite{batzner_efficientad_2024,akcay_anomalib_2022} In this case, the detector acts as a region proposal mechanism, while the downstream anomaly model estimates how strongly each loaf deviates from the normal appearance learned from production data.\cite{labby_thesis_repo_2026}
